\section{Introduction}
Quantum dot sensitized solar cells (QDSSCs) have attracted significant scientific interest as a promising third-generation photovoltaic technology. Their low manufacturing cost, ease of fabrication, and high theoretical limit of up to $44\%$ via multiple exciton generation make them viable candidates for large-scale energy harvesting. 

In a typical QDSSC, cadmium chalcogenide quantum dots (e.g., CdS, CdSe) serve as light harvesters sensitizing a wide-bandgap metal oxide electron transport layer (ETL), most commonly titanium dioxide ($\text{TiO}_2$). However, conventional random-nanoparticle $\text{TiO}_2$ networks suffer from slow electron mobility and high rates of interfacial charge recombination. Photoexcited electrons face numerous grain boundaries and defects during diffusion, leading to poor collection efficiency.

To address these limitations, we introduce the concept of ``synapse engineering'' to the solar cell interface. Similar to biological synapses that optimize signal transmission via nanoscale junctions, synapse-engineered ETLs utilize highly ordered, passivated nano-junctions to facilitate fast charge collection. In this work, we present a hydrothermal synthesis route using a specialized block-copolymer surfactant, Apex-9, to assemble highly ordered mesoporous $\text{TiO}_2$ structures. This structure dramatically improves quantum dot loading and curtails back-electron transfer.
